\documentclass[11pt,a4paper]{article}

%========================
% PACKAGES
%========================
\usepackage[utf8]{inputenc}
\usepackage[T1]{fontenc}
\usepackage[french]{babel}
\usepackage{lmodern}
\usepackage{geometry}
\usepackage{amsmath,amssymb,amsthm}
\usepackage{siunitx}
\usepackage{physics}
\usepackage{enumitem}
\usepackage{fancyhdr}
\usepackage{xcolor}
\usepackage[most]{tcolorbox}
\usepackage{multicol}
\usepackage{array}
\usepackage{booktabs}
\usepackage{hyperref}

\geometry{margin=2cm}
\setlength{\parindent}{0pt}
\setlength{\parskip}{0.45em}
\sisetup{
  locale = FR,
  per-mode = symbol,
  detect-all
}

%========================
% COULEURS ET BOÎTES
%========================
\definecolor{bleuprepa}{RGB}{20,55,110}
\definecolor{vertprepa}{RGB}{20,105,70}
\definecolor{rougeprepa}{RGB}{160,35,35}
\definecolor{orangeprepa}{RGB}{190,100,20}
\definecolor{grisclair}{RGB}{245,247,250}

\tcbset{
  enhanced,
  breakable,
  sharp corners,
  boxrule=0.7pt,
  colback=white,
  fonttitle=\bfseries
}

\newtcolorbox{methodebox}[1]{
  colframe=bleuprepa,
  colback=blue!3!white,
  title={Méthode -- #1}
}

\newtcolorbox{resultatbox}{
  colframe=vertprepa,
  colback=green!3!white,
  title={Résultat}
}

\newtcolorbox{remarquebox}{
  colframe=orangeprepa,
  colback=orange!4!white,
  title={Remarque physique}
}

\newtcolorbox{attentionbox}{
  colframe=rougeprepa,
  colback=red!3!white,
  title={Point de vigilance}
}

%========================
% EN-TÊTE
%========================
\pagestyle{fancy}
\fancyhf{}
\lhead{\textbf{Terminale générale -- Spécialité Physique-Chimie}}
\rhead{\textbf{Corrigé détaillé -- Mécanique}}
\cfoot{\thepage}

\title{
  \vspace{-1cm}
  \Huge\textbf{Corrigé complet et détaillé}\\
  \Large Devoir surveillé -- Mécanique\\
  \large Niveau Terminale approfondi, esprit prépa
}
\author{Chapitre : Mouvement et interactions}
\date{}

\begin{document}

\maketitle

\begin{center}
\begin{tcolorbox}[colframe=bleuprepa,colback=grisclair,width=0.95\textwidth]
Ce corrigé insiste sur la méthode : choix du système, référentiel, bilan des forces,
application de la deuxième loi de Newton, projection, intégration, puis interprétation physique.
Les résultats numériques sont donnés avec des unités cohérentes.
\end{tcolorbox}
\end{center}

\tableofcontents
\newpage

%========================================================
\section{Exercice 1 -- Sécurisation d'une piste de saut à ski nautique}

On étudie le mouvement du centre d'inertie $G$ d'un skieur après son départ de la rampe.
Le repère est $(O;\vec e_x,\vec e_z)$, avec $\vec e_x$ horizontal dans le sens du mouvement et
$\vec e_z$ vertical vers le haut. Le point d'impact avec l'eau correspond à $z=0$.

Données :
\[
h=\SI{2.40}{m}, \qquad v_0=\SI{54}{km.h^{-1}}, \qquad g=\SI{9.81}{m.s^{-2}}.
\]

\subsection{Partie A -- Mise en équation du mouvement}

\subsubsection*{A.1. Système et référentiel}

Le système étudié est le \textbf{skieur équipé}, modélisé par son centre d'inertie $G$.
Le référentiel choisi est le \textbf{référentiel terrestre}, supposé galiléen pendant la durée courte du saut.

\begin{resultatbox}
Système : skieur équipé. \\
Référentiel : référentiel terrestre supposé galiléen.
\end{resultatbox}

\subsubsection*{A.2. Bilan des forces}

Après le départ de la rampe, et dans cette première modélisation, on néglige les frottements de l'air.
Le skieur n'est plus en contact avec la rampe. La seule force exercée sur lui est donc son poids :
\[
\vec P = m\vec g = -mg\,\vec e_z.
\]

\begin{attentionbox}
Il ne faut pas ajouter une force ``de vitesse'' ou une force ``du mouvement''. En mécanique newtonienne,
la vitesse n'est pas une force. Une fois le skieur lâché, il conserve une vitesse horizontale par inertie.
\end{attentionbox}

\subsubsection*{A.3. Deuxième loi de Newton et accélération}

D'après la deuxième loi de Newton :
\[
\sum \vec F_{\text{ext}} = m\vec a_G.
\]
Ici :
\[
\vec P = m\vec a_G.
\]
Donc :
\[
-mg\,\vec e_z = m\vec a_G.
\]
En divisant par $m$ :
\[
\vec a_G = -g\,\vec e_z.
\]
Ainsi, dans la base $(\vec e_x,\vec e_z)$ :
\[
a_x=0,\qquad a_z=-g.
\]

\begin{resultatbox}
\[
\boxed{\vec a_G = -g\,\vec e_z}
\qquad\text{c'est-à-dire}\qquad
\boxed{a_x=0,\quad a_z=-g.}
\]
\end{resultatbox}

\subsubsection*{A.4. Équations horaires}

La vitesse initiale est horizontale :
\[
\vec v(0)=v_0\vec e_x.
\]
Donc :
\[
v_x(0)=v_0,\qquad v_z(0)=0.
\]

Puisque $a_x=0$ :
\[
\frac{dv_x}{dt}=0 \quad\Rightarrow\quad v_x(t)=v_0.
\]
Puisque $a_z=-g$ :
\[
\frac{dv_z}{dt}=-g \quad\Rightarrow\quad v_z(t)=-gt,
\]
car $v_z(0)=0$.

On intègre encore pour obtenir les positions. Le point $O$ est pris à la sortie de la rampe et l'eau est à $z=0$.
Le skieur quitte la rampe à la hauteur $h$ au-dessus de l'eau. On prend donc :
\[
x(0)=0,\qquad z(0)=h.
\]
Ainsi :
\[
x(t)=v_0t,
\]
et
\[
z(t)=h-\frac12gt^2.
\]

\begin{resultatbox}
\[
\boxed{x(t)=v_0t}
\qquad\text{et}\qquad
\boxed{z(t)=h-\frac12gt^2.}
\]
\end{resultatbox}

\subsubsection*{A.5. Équation cartésienne de la trajectoire}

À partir de :
\[
x(t)=v_0t,
\]
on obtient :
\[
t=\frac{x}{v_0}.
\]
On remplace dans $z(t)$ :
\[
z=h-\frac12g\left(\frac{x}{v_0}\right)^2.
\]
Donc :
\[
z=h-\frac{g}{2v_0^2}x^2.
\]

\begin{resultatbox}
\[
\boxed{z(x)=h-\frac{g}{2v_0^2}x^2.}
\]
La trajectoire est une parabole tournée vers le bas.
\end{resultatbox}

\begin{remarquebox}
L'absence de force horizontale explique que le mouvement horizontal soit uniforme. Le mouvement vertical est uniformément accéléré vers le bas.
Le mouvement total est donc la composition d'un mouvement rectiligne uniforme horizontal et d'une chute libre verticale.
\end{remarquebox}

\subsection{Partie B -- Exploitation numérique}

\subsubsection*{B.1. Conversion de la vitesse initiale}

\[
v_0=\SI{54}{km.h^{-1}}
=54\times \frac{1000}{3600}\ \si{m.s^{-1}}
=15.0\ \si{m.s^{-1}}.
\]

\begin{resultatbox}
\[
\boxed{v_0=\SI{15.0}{m.s^{-1}}.}
\]
\end{resultatbox}

\subsubsection*{B.2. Durée du vol}

Le contact avec l'eau a lieu lorsque $z(t_f)=0$ :
\[
h-\frac12gt_f^2=0.
\]
Donc :
\[
\frac12gt_f^2=h,
\]
d'où :
\[
t_f^2=\frac{2h}{g}.
\]
Ainsi :
\[
t_f=\sqrt{\frac{2h}{g}}.
\]
Numériquement :
\[
t_f=\sqrt{\frac{2\times 2.40}{9.81}}
=\sqrt{0.489}
\simeq 0.699\ \si{s}.
\]

\begin{resultatbox}
\[
\boxed{t_f\simeq \SI{0.70}{s}.}
\]
\end{resultatbox}

\subsubsection*{B.3. Portée horizontale}

La portée horizontale vaut :
\[
D=x(t_f)=v_0t_f.
\]
Donc :
\[
D=15.0\times 0.699\simeq 10.5\ \si{m}.
\]

\begin{resultatbox}
\[
\boxed{D\simeq \SI{10.5}{m}.}
\]
\end{resultatbox}

\subsubsection*{B.4. Vecteur vitesse à l'impact}

On a :
\[
v_x(t)=v_0=15.0\ \si{m.s^{-1}},
\]
et :
\[
v_z(t)=-gt.
\]
À l'impact :
\[
v_z(t_f)=-9.81\times 0.699\simeq -6.86\ \si{m.s^{-1}}.
\]
Donc :
\[
\vec v(t_f)=15.0\,\vec e_x-6.86\,\vec e_z.
\]

\begin{resultatbox}
\[
\boxed{\vec v(t_f)\simeq 15.0\,\vec e_x-6.86\,\vec e_z\quad (\si{m.s^{-1}}).}
\]
\end{resultatbox}

\subsubsection*{B.5. Valeur de la vitesse à l'impact}

La norme vaut :
\[
v_f=\sqrt{v_x^2+v_z^2}.
\]
Donc :
\[
v_f=\sqrt{15.0^2+6.86^2}
=\sqrt{225+47.1}
=\sqrt{272.1}
\simeq 16.5\ \si{m.s^{-1}}.
\]
En km/h :
\[
16.5\times 3.6\simeq 59.4\ \si{km.h^{-1}}.
\]

\begin{resultatbox}
\[
\boxed{v_f\simeq \SI{16.5}{m.s^{-1}}\simeq \SI{59.4}{km.h^{-1}}.}
\]
\end{resultatbox}

\begin{remarquebox}
La vitesse à l'impact est supérieure à la vitesse initiale car le skieur a perdu de l'énergie potentielle de pesanteur,
transformée en énergie cinétique. La composante horizontale reste constante, mais une composante verticale descendante apparaît.
\end{remarquebox}

\subsection{Partie C -- Sécurité du bassin de réception}

Le bassin commence à $\SI{5.0}{m}$ après la rampe et s'étend sur $\SI{18.0}{m}$.
Il couvre donc l'intervalle :
\[
[5.0;23.0]\ \si{m}.
\]
On a trouvé :
\[
D\simeq \SI{10.5}{m}.
\]

\subsubsection*{C.1. Le skieur tombe-t-il dans la zone prévue ?}

Comme :
\[
5.0 < 10.5 < 23.0,
\]
le skieur tombe dans le bassin.

\begin{resultatbox}
\[
\boxed{\text{Oui, le skieur tombe dans la zone prévue.}}
\]
\end{resultatbox}

\subsubsection*{C.2. La vitesse de \SI{54}{km.h^{-1}} est-elle adaptée pour viser entre 8,0 m et 15,0 m ?}

L'intervalle souhaité est :
\[
[8.0;15.0]\ \si{m}.
\]
Or :
\[
D\simeq \SI{10.5}{m}.
\]
Donc :
\[
8.0<10.5<15.0.
\]

\begin{resultatbox}
\[
\boxed{\text{Oui, la vitesse initiale de \SI{54}{km.h^{-1}} est adaptée dans ce modèle.}}
\]
\end{resultatbox}

\subsubsection*{C.3. Vitesse minimale pour atteindre \SI{8.0}{m}}

Le temps de chute dépend uniquement du mouvement vertical :
\[
t_f=\sqrt{\frac{2h}{g}}.
\]
Pour atteindre au moins $D_{\min}=\SI{8.0}{m}$ :
\[
D_{\min}=v_{0,\min}t_f.
\]
Donc :
\[
v_{0,\min}=\frac{D_{\min}}{t_f}.
\]
Numériquement :
\[
v_{0,\min}=\frac{8.0}{0.699}\simeq 11.4\ \si{m.s^{-1}}.
\]
En km/h :
\[
v_{0,\min}\simeq 11.4\times 3.6\simeq 41.2\ \si{km.h^{-1}}.
\]

\begin{resultatbox}
\[
\boxed{v_{0,\min}\simeq \SI{11.4}{m.s^{-1}}\simeq \SI{41.2}{km.h^{-1}}.}
\]
\end{resultatbox}

\subsection{Partie D -- Prise en compte simplifiée de l'air}

On ajoute une force horizontale constante :
\[
\vec F_{\text{air}}=-f\vec e_x,
\qquad f=\SI{60}{N},
\qquad m=\SI{75}{kg}.
\]

\subsubsection*{D.1. Nouvelles coordonnées de l'accélération}

Les forces sont :
\[
\vec P=-mg\vec e_z,
\qquad
\vec F_{\text{air}}=-f\vec e_x.
\]
La deuxième loi de Newton donne :
\[
m\vec a=-f\vec e_x-mg\vec e_z.
\]
Donc :
\[
\vec a=-\frac{f}{m}\vec e_x-g\vec e_z.
\]
Ainsi :
\[
a_x=-\frac{f}{m}=-\frac{60}{75}=-0.80\ \si{m.s^{-2}},
\]
et :
\[
a_z=-g=-9.81\ \si{m.s^{-2}}.
\]

\begin{resultatbox}
\[
\boxed{a_x=-0.80\ \si{m.s^{-2}},\qquad a_z=-9.81\ \si{m.s^{-2}}.}
\]
\end{resultatbox}

\subsubsection*{D.2. Nouvelle expression de $x(t)$}

On a :
\[
\frac{dv_x}{dt}=-\frac{f}{m}.
\]
Donc :
\[
v_x(t)=v_0-\frac{f}{m}t.
\]
Puis :
\[
x(t)=v_0t-\frac12\frac{f}{m}t^2.
\]

\begin{resultatbox}
\[
\boxed{x(t)=v_0t-\frac{f}{2m}t^2.}
\]
\end{resultatbox}

\subsubsection*{D.3. Pourquoi le temps de chute reste-t-il inchangé ?}

La force de l'air introduite est supposée uniquement horizontale :
\[
\vec F_{\text{air}}=-f\vec e_x.
\]
Elle ne possède donc aucune composante selon $\vec e_z$. L'équation verticale reste :
\[
z(t)=h-\frac12gt^2.
\]
Le temps de chute, déterminé par $z(t_f)=0$, reste donc :
\[
t_f=\sqrt{\frac{2h}{g}}.
\]

\begin{resultatbox}
\[
\boxed{\text{Le temps de chute reste inchangé car le mouvement vertical n'est pas modifié.}}
\]
\end{resultatbox}

\subsubsection*{D.4. Nouvelle portée horizontale}

On utilise le même temps :
\[
t_f\simeq 0.699\ \si{s}.
\]
La nouvelle portée vaut :
\[
D'=v_0t_f-\frac{f}{2m}t_f^2.
\]
Or :
\[
\frac{f}{2m}=\frac{60}{2\times 75}=0.40.
\]
Donc :
\[
D'=15.0\times 0.699-0.40\times (0.699)^2.
\]
\[
D'\simeq 10.49-0.196\simeq 10.29\ \si{m}.
\]

\begin{resultatbox}
\[
\boxed{D'\simeq \SI{10.3}{m}.}
\]
\end{resultatbox}

\subsubsection*{D.5. Effet sur la sécurité du saut}

La force de l'air ralentit le mouvement horizontal : la portée diminue.
Ici, la diminution est faible :
\[
10.5\ \si{m}\quad\longrightarrow\quad 10.3\ \si{m}.
\]
Le skieur reste dans la zone souhaitée $[8.0;15.0]\ \si{m}$.

\begin{resultatbox}
\[
\boxed{\text{Dans ce modèle, l'air réduit légèrement la portée, sans compromettre la sécurité.}}
\]
\end{resultatbox}

%========================================================
\newpage
\section{Exercice 2 -- Virage d'une voiture sur route circulaire}

Données :
\[
m=1.20\times 10^3\ \si{kg},\qquad R=\SI{45}{m},\qquad \mu=0.65.
\]

\subsection{Partie A -- Analyse qualitative}

\subsubsection*{A.1. Mouvement uniforme ? Mouvement accéléré ?}

La voiture roule à vitesse de valeur constante. Le mouvement est donc \textbf{uniforme} au sens où la norme de la vitesse est constante.

Mais la trajectoire est circulaire : la direction du vecteur vitesse change constamment. Or l'accélération mesure la variation du vecteur vitesse, pas seulement la variation de sa norme.

Donc le mouvement est \textbf{accéléré}, avec une accélération centripète dirigée vers le centre du cercle.

\begin{resultatbox}
\[
\boxed{\text{Le mouvement est uniforme mais accéléré.}}
\]
\end{resultatbox}

\subsubsection*{A.2. Représentation de $\vec v$ et $\vec a$}

En un point du virage :
\[
\vec v
\]
est tangent à la trajectoire, dans le sens du mouvement.

\[
\vec a
\]
est radial, dirigé vers le centre du cercle.

\begin{methodebox}{Schéma attendu}
Sur un cercle, placer la voiture en un point. Tracer $\vec v$ tangent à la trajectoire et $\vec a$ perpendiculaire à $\vec v$, orienté vers le centre du cercle.
\end{methodebox}

\subsubsection*{A.3. Valeur de l'accélération}

Dans un mouvement circulaire uniforme :
\[
a=\frac{v^2}{R}.
\]

\begin{resultatbox}
\[
\boxed{a=\frac{v^2}{R}.}
\]
\end{resultatbox}

\subsection{Partie B -- Forces et condition dynamique}

\subsubsection*{B.1. Bilan des forces}

Sur une route horizontale, les forces extérieures exercées sur la voiture sont :
\begin{itemize}
  \item le poids : $\vec P=m\vec g$, vertical vers le bas ;
  \item la réaction normale de la route : $\vec N$, verticale vers le haut ;
  \item la force d'adhérence latérale $\vec F_{\text{adh}}$, horizontale, dirigée vers le centre du virage.
\end{itemize}

\subsubsection*{B.2. Résultante verticale}

La voiture ne décolle pas et ne s'enfonce pas dans la route. Son mouvement vertical est nul.
Donc :
\[
a_z=0.
\]
D'après Newton :
\[
N-mg=0.
\]
Donc :
\[
N=mg.
\]

\begin{resultatbox}
\[
\boxed{N=mg.}
\]
\end{resultatbox}

\subsubsection*{B.3. Force responsable de la courbure}

La force qui permet de courber la trajectoire est la force d'adhérence latérale.
Elle fournit la résultante centripète.

\begin{resultatbox}
\[
\boxed{\vec F_{\text{adh}} \text{ est responsable de la courbure de la trajectoire.}}
\]
\end{resultatbox}

\subsubsection*{B.4. Relation entre $F_{\text{adh}}$ et $v$}

On projette la deuxième loi de Newton selon l'axe radial, orienté vers le centre :
\[
\sum F_r=ma_r.
\]
La seule force horizontale radiale est l'adhérence :
\[
F_{\text{adh}}=m\frac{v^2}{R}.
\]

\begin{resultatbox}
\[
\boxed{F_{\text{adh}}=m\frac{v^2}{R}.}
\]
\end{resultatbox}

\begin{remarquebox}
Plus la vitesse est grande, plus la force latérale nécessaire augmente comme $v^2$.
Doubler la vitesse nécessite donc quatre fois plus d'adhérence.
\end{remarquebox}

\subsection{Partie C -- Vitesse maximale sans glisser}

\subsubsection*{C.1. Vitesse maximale}

La valeur maximale de la force d'adhérence est :
\[
F_{\max}=\mu mg.
\]
Pour ne pas glisser :
\[
F_{\text{adh}}\leq F_{\max}.
\]
Donc :
\[
m\frac{v^2}{R}\leq \mu mg.
\]
On simplifie par $m$ :
\[
\frac{v^2}{R}\leq \mu g.
\]
Donc :
\[
v^2\leq \mu gR.
\]
La vitesse maximale vaut :
\[
v_{\max}=\sqrt{\mu gR}.
\]
Numériquement :
\[
v_{\max}=\sqrt{0.65\times 9.81\times 45}
=\sqrt{286.94}
\simeq 16.9\ \si{m.s^{-1}}.
\]

\begin{resultatbox}
\[
\boxed{v_{\max}\simeq \SI{16.9}{m.s^{-1}}.}
\]
\end{resultatbox}

\subsubsection*{C.2. Conversion en km/h}

\[
v_{\max}\simeq 16.9\times 3.6\simeq 61.0\ \si{km.h^{-1}}.
\]

\begin{resultatbox}
\[
\boxed{v_{\max}\simeq \SI{61}{km.h^{-1}}.}
\]
\end{resultatbox}

\subsubsection*{C.3. Voiture à \SI{80}{km.h^{-1}}}

On convertit :
\[
80\ \si{km.h^{-1}}=\frac{80}{3.6}\simeq 22.2\ \si{m.s^{-1}}.
\]
Or :
\[
22.2>16.9.
\]
Donc la force d'adhérence nécessaire dépasse la force maximale disponible.

\begin{resultatbox}
\[
\boxed{\text{À \SI{80}{km.h^{-1}}, le virage n'est pas théoriquement franchissable sans glisser.}}
\]
\end{resultatbox}

\subsection{Partie D -- Influence de la pluie}

Sous la pluie :
\[
\mu'=0.35.
\]

\subsubsection*{D.1. Nouvelle vitesse maximale}

\[
v'_{\max}=\sqrt{\mu' gR}.
\]
Donc :
\[
v'_{\max}=\sqrt{0.35\times 9.81\times 45}
=\sqrt{154.51}
\simeq 12.4\ \si{m.s^{-1}}.
\]
En km/h :
\[
v'_{\max}\simeq 12.4\times 3.6\simeq 44.7\ \si{km.h^{-1}}.
\]

\begin{resultatbox}
\[
\boxed{v'_{\max}\simeq \SI{12.4}{m.s^{-1}}\simeq \SI{44.7}{km.h^{-1}}.}
\]
\end{resultatbox}

\subsubsection*{D.2. Comparaison et risque routier}

Sur route sèche :
\[
v_{\max}\simeq \SI{61}{km.h^{-1}}.
\]
Sous la pluie :
\[
v'_{\max}\simeq \SI{45}{km.h^{-1}}.
\]
La vitesse maximale baisse fortement.

\begin{remarquebox}
La pluie diminue le coefficient d'adhérence. Le conducteur doit réduire sa vitesse, car la force nécessaire pour tourner varie comme $v^2$.
Une vitesse raisonnable sur route sèche peut devenir dangereuse sous la pluie.
\end{remarquebox}

\subsubsection*{D.3. Rapport des vitesses maximales}

On a :
\[
v_{\max}=\sqrt{\mu gR},
\qquad
v'_{\max}=\sqrt{\mu' gR}.
\]
Donc :
\[
\frac{v'_{\max}}{v_{\max}}
=
\frac{\sqrt{\mu'gR}}{\sqrt{\mu gR}}
=
\sqrt{\frac{\mu'gR}{\mu gR}}
=
\sqrt{\frac{\mu'}{\mu}}.
\]

\begin{resultatbox}
\[
\boxed{\frac{v'_{\max}}{v_{\max}}=\sqrt{\frac{\mu'}{\mu}}.}
\]
\end{resultatbox}

\subsection{Partie E -- Question de modélisation}

\subsubsection*{E.1. Pourquoi la phrase du conducteur est fausse ?}

Le conducteur affirme :
\[
\text{``Si je garde la même valeur de vitesse, je n'accélère pas.''}
\]
Cette phrase confond la valeur de la vitesse et le vecteur vitesse.

En mécanique, l'accélération est :
\[
\vec a=\frac{d\vec v}{dt}.
\]
Même si la norme $v=\norm{\vec v}$ est constante, le vecteur $\vec v$ peut varier si sa direction change.
Dans un virage, la direction de $\vec v$ change sans cesse.

\begin{resultatbox}
\[
\boxed{\text{La phrase est fausse : la vitesse vectorielle varie car sa direction change.}}
\]
\end{resultatbox}

\subsubsection*{E.2. Que mesure l'accélération ?}

L'accélération mesure la variation du vecteur vitesse :
\[
\vec a=\frac{d\vec v}{dt}.
\]
Elle peut traduire :
\begin{itemize}
  \item une variation de la valeur de la vitesse ;
  \item une variation de la direction de la vitesse ;
  \item ou les deux.
\end{itemize}
Dans un mouvement circulaire uniforme, elle mesure uniquement la variation de direction de $\vec v$.

\begin{resultatbox}
\[
\boxed{\text{Ici, l'accélération mesure la variation de direction du vecteur vitesse.}}
\]
\end{resultatbox}

%========================================================
\newpage
\section{Exercice 3 -- Comment viser le plus loin possible ?}

Données :
\[
v_0=\SI{18.0}{m.s^{-1}},\qquad g=\SI{9.81}{m.s^{-2}}.
\]
Le tir part du sol depuis $O$ :
\[
x(0)=0,\qquad z(0)=0.
\]

\subsection{Partie A -- Équations horaires}

\subsubsection*{A.1. Coordonnées initiales de la vitesse}

La vitesse initiale fait un angle $\alpha$ avec l'horizontale :
\[
\vec v_0=v_0\cos\alpha\,\vec e_x+v_0\sin\alpha\,\vec e_z.
\]

\begin{resultatbox}
\[
\boxed{v_{0x}=v_0\cos\alpha,\qquad v_{0z}=v_0\sin\alpha.}
\]
\end{resultatbox}

\subsubsection*{A.2. Équations horaires}

Sans frottements, seule la pesanteur agit :
\[
\vec a=-g\vec e_z.
\]
Donc :
\[
a_x=0,\qquad a_z=-g.
\]
En intégrant :
\[
v_x(t)=v_0\cos\alpha,
\qquad
v_z(t)=v_0\sin\alpha-gt.
\]
Puis :
\[
x(t)=v_0\cos\alpha\,t,
\]
et :
\[
z(t)=v_0\sin\alpha\,t-\frac12gt^2.
\]

\begin{resultatbox}
\[
\boxed{x(t)=v_0\cos\alpha\,t}
\qquad
\boxed{z(t)=v_0\sin\alpha\,t-\frac12gt^2.}
\]
\end{resultatbox}

\subsubsection*{A.3. Équation cartésienne}

Si $\cos\alpha\neq 0$, on a :
\[
t=\frac{x}{v_0\cos\alpha}.
\]
Donc :
\[
z=x\tan\alpha-\frac12g\left(\frac{x}{v_0\cos\alpha}\right)^2.
\]
Ainsi :
\[
z=x\tan\alpha-\frac{g}{2v_0^2\cos^2\alpha}x^2.
\]

\begin{resultatbox}
\[
\boxed{z(x)=x\tan\alpha-\frac{g}{2v_0^2\cos^2\alpha}x^2.}
\]
\end{resultatbox}

\subsection{Partie B -- Temps de vol et portée}

\subsubsection*{B.1. Instant non nul du retour au sol}

Le retour au sol correspond à :
\[
z(t)=0.
\]
Donc :
\[
v_0\sin\alpha\,t-\frac12gt^2=0.
\]
On factorise :
\[
t\left(v_0\sin\alpha-\frac12gt\right)=0.
\]
Les solutions sont :
\[
t=0
\]
et
\[
v_0\sin\alpha-\frac12gt=0.
\]
Donc :
\[
t=\frac{2v_0\sin\alpha}{g}.
\]

\begin{resultatbox}
\[
\boxed{t_{\text{vol}}=\frac{2v_0\sin\alpha}{g}.}
\]
\end{resultatbox}

\subsubsection*{B.2. Portée horizontale}

La portée vaut :
\[
D(\alpha)=x(t_{\text{vol}}).
\]
Donc :
\[
D(\alpha)=v_0\cos\alpha \times \frac{2v_0\sin\alpha}{g}.
\]
Ainsi :
\[
D(\alpha)=\frac{2v_0^2\sin\alpha\cos\alpha}{g}.
\]
Or :
\[
\sin(2\alpha)=2\sin\alpha\cos\alpha.
\]
Donc :
\[
D(\alpha)=\frac{v_0^2}{g}\sin(2\alpha).
\]

\begin{resultatbox}
\[
\boxed{D(\alpha)=\frac{v_0^2}{g}\sin(2\alpha).}
\]
\end{resultatbox}

\subsubsection*{B.3. Homogénéité}

\[
[v_0^2]=\si{m^2.s^{-2}},
\qquad
[g]=\si{m.s^{-2}}.
\]
Donc :
\[
\left[\frac{v_0^2}{g}\right]
=
\frac{\si{m^2.s^{-2}}}{\si{m.s^{-2}}}
=
\si{m}.
\]
Comme $\sin(2\alpha)$ est sans dimension :
\[
[D]=\si{m}.
\]

\begin{resultatbox}
\[
\boxed{\text{L'expression est homogène à une longueur.}}
\]
\end{resultatbox}

\subsection{Partie C -- Angle de portée maximale}

\subsubsection*{C.1. Condition sur $\sin(2\alpha)$}

Comme :
\[
D(\alpha)=\frac{v_0^2}{g}\sin(2\alpha),
\]
et que $\frac{v_0^2}{g}>0$, la portée est maximale lorsque :
\[
\sin(2\alpha)=1.
\]

\subsubsection*{C.2. Angle maximal}

\[
\sin(2\alpha)=1
\quad\Longleftrightarrow\quad
2\alpha=90^\circ
\]
dans le domaine usuel $0^\circ<\alpha<90^\circ$.
Donc :
\[
\alpha=45^\circ.
\]

\begin{resultatbox}
\[
\boxed{\alpha_{\max}=45^\circ.}
\]
\end{resultatbox}

\subsubsection*{C.3. Portée maximale}

\[
D_{\max}=\frac{v_0^2}{g}.
\]
Numériquement :
\[
D_{\max}=\frac{18.0^2}{9.81}
=\frac{324}{9.81}
\simeq 33.0\ \si{m}.
\]

\begin{resultatbox}
\[
\boxed{D_{\max}\simeq \SI{33.0}{m}.}
\]
\end{resultatbox}

\subsection{Partie D -- Symétrie des tirs complémentaires}

\subsubsection*{D.1. Même portée pour $\alpha+\beta=90^\circ$}

Si :
\[
\alpha+\beta=90^\circ,
\]
alors :
\[
\beta=90^\circ-\alpha.
\]
Donc :
\[
2\beta=180^\circ-2\alpha.
\]
Or :
\[
\sin(180^\circ-\theta)=\sin\theta.
\]
Ainsi :
\[
\sin(2\beta)=\sin(2\alpha).
\]
Donc :
\[
D(\beta)=\frac{v_0^2}{g}\sin(2\beta)
=
\frac{v_0^2}{g}\sin(2\alpha)
=
D(\alpha).
\]

\begin{resultatbox}
\[
\boxed{D(\alpha)=D(\beta)\quad\text{si}\quad \alpha+\beta=90^\circ.}
\]
\end{resultatbox}

\subsubsection*{D.2. Interprétation physique}

Deux tirs complémentaires donnent la même portée :
\begin{itemize}
  \item l'un est plus tendu : grande vitesse horizontale, temps de vol court ;
  \item l'autre est plus en cloche : vitesse horizontale plus faible, temps de vol plus long.
\end{itemize}
Les deux effets se compensent exactement dans le modèle sans frottements et avec départ et arrivée à la même hauteur.

\subsubsection*{D.3. Angles pour atteindre \SI{20.0}{m}}

On cherche :
\[
D(\alpha)=20.0.
\]
Donc :
\[
\frac{v_0^2}{g}\sin(2\alpha)=20.0.
\]
Ainsi :
\[
\sin(2\alpha)=\frac{20.0\,g}{v_0^2}.
\]
Numériquement :
\[
\sin(2\alpha)=\frac{20.0\times 9.81}{18.0^2}
=\frac{196.2}{324}
\simeq 0.606.
\]
On cherche les solutions dans $0^\circ<\alpha<90^\circ$.

\[
2\alpha_1=\arcsin(0.606)\simeq 37.3^\circ.
\]
Donc :
\[
\alpha_1\simeq 18.7^\circ.
\]
L'autre solution est :
\[
2\alpha_2=180^\circ-37.3^\circ=142.7^\circ.
\]
Donc :
\[
\alpha_2\simeq 71.3^\circ.
\]

\begin{resultatbox}
\[
\boxed{\alpha_1\simeq 18.7^\circ,\qquad \alpha_2\simeq 71.3^\circ.}
\]
\end{resultatbox}

\subsection{Partie E -- Atteindre une cible en hauteur}

Données :
\[
x_c=\SI{12.0}{m},\qquad z_c=\SI{3.00}{m}.
\]

\subsubsection*{E.1. Équation vérifiée par $\alpha$}

La trajectoire est :
\[
z(x)=x\tan\alpha-\frac{g}{2v_0^2\cos^2\alpha}x^2.
\]
Pour passer par $(x_c,z_c)$ :
\[
z_c=x_c\tan\alpha-\frac{g}{2v_0^2\cos^2\alpha}x_c^2.
\]

\begin{resultatbox}
\[
\boxed{
z_c=x_c\tan\alpha-\frac{g x_c^2}{2v_0^2\cos^2\alpha}.
}
\]
\end{resultatbox}

On peut aussi poser :
\[
u=\tan\alpha.
\]
Comme :
\[
\frac{1}{\cos^2\alpha}=1+\tan^2\alpha=1+u^2,
\]
l'équation devient :
\[
z_c=x_cu-\frac{gx_c^2}{2v_0^2}(1+u^2).
\]
C'est une équation du second degré en $u=\tan\alpha$.

\subsubsection*{E.2. Pourquoi zéro, une ou deux trajectoires ?}

L'équation en $u=\tan\alpha$ est une équation du second degré. Selon son discriminant :
\[
\Delta>0 : \text{deux solutions réelles},
\]
\[
\Delta=0 : \text{une solution réelle double},
\]
\[
\Delta<0 : \text{aucune solution réelle}.
\]
Physiquement :
\begin{itemize}
  \item deux solutions : un tir tendu et un tir en cloche ;
  \item une solution : la cible est atteinte de manière limite ;
  \item zéro solution : la cible est trop loin ou trop haute pour la vitesse initiale donnée.
\end{itemize}

\subsubsection*{E.3. Interprétation des deux solutions éventuelles}

S'il y a deux angles :
\begin{itemize}
  \item le plus petit angle correspond à une trajectoire tendue, rapide, avec faible temps de vol ;
  \item le plus grand angle correspond à une trajectoire plus haute, plus longue en temps, dite en cloche.
\end{itemize}

\begin{remarquebox}
Cette dualité est très générale dans les tirs paraboliques : pour atteindre une cible accessible,
il existe souvent deux arcs possibles, sauf au point limite où les deux se confondent.
\end{remarquebox}

%========================================================
\newpage
\section{Exercice 4 -- Force de recentrage dynamique}

On introduit une force :
\[
\vec F_r=-kx\,\vec e_x,
\]
où $k>0$.

Un mobile de masse $m$ se déplace sur un rail horizontal sans frottement.

\subsection{Partie A -- Compréhension de la force}

\subsubsection*{A.1. Dimension physique de $k$}

On a :
\[
F=kx.
\]
Donc :
\[
k=\frac{F}{x}.
\]
Or :
\[
[F]=\si{N},\qquad [x]=\si{m}.
\]
Donc :
\[
[k]=\si{N.m^{-1}}.
\]

\begin{resultatbox}
\[
\boxed{[k]=\si{N.m^{-1}}.}
\]
\end{resultatbox}

En unités fondamentales :
\[
\si{N.m^{-1}}=\frac{\si{kg.m.s^{-2}}}{\si{m}}
=\si{kg.s^{-2}}.
\]

\subsubsection*{A.2. Sens de la force}

\[
\vec F_r=-kx\vec e_x.
\]

Si $x>0$, alors $-kx<0$, donc :
\[
\vec F_r \text{ est dirigée vers } -\vec e_x.
\]
Elle pointe vers la gauche, donc vers l'origine.

Si $x<0$, alors $-kx>0$, donc :
\[
\vec F_r \text{ est dirigée vers } +\vec e_x.
\]
Elle pointe vers la droite, donc encore vers l'origine.

\begin{resultatbox}
\[
\boxed{
x>0 : \vec F_r \text{ vers la gauche};\qquad
x<0 : \vec F_r \text{ vers la droite}.
}
\]
\end{resultatbox}

\subsubsection*{A.3. Valeur de la force à l'origine}

Si $x=0$ :
\[
\vec F_r=-k\times 0\,\vec e_x=\vec 0.
\]

\begin{resultatbox}
\[
\boxed{\vec F_r=\vec 0 \quad \text{pour } x=0.}
\]
\end{resultatbox}

\subsubsection*{A.4. Pourquoi le nom est pertinent ?}

La force est toujours orientée vers l'origine. Elle tend donc à ramener le mobile vers $x=0$.
Le nom ``force de recentrage'' est pertinent car $x=0$ joue le rôle de position d'équilibre.

\begin{resultatbox}
\[
\boxed{\text{La force recentre le mobile vers l'origine.}}
\]
\end{resultatbox}

\subsection{Partie B -- Conséquence dynamique}

\subsubsection*{B.1. Deuxième loi de Newton}

Sur l'axe horizontal, la seule force est :
\[
F_x=-kx.
\]
D'après Newton :
\[
ma_x=-kx.
\]
Donc :
\[
a_x=-\frac{k}{m}x.
\]

\begin{resultatbox}
\[
\boxed{a_x=-\frac{k}{m}x.}
\]
\end{resultatbox}

\subsubsection*{B.2. Cas $x>0$}

Si $x>0$, alors :
\[
a_x=-\frac{k}{m}x<0.
\]
L'accélération est dirigée vers la gauche, donc vers l'origine.

\begin{resultatbox}
\[
\boxed{x>0 \Longrightarrow a_x<0.}
\]
\end{resultatbox}

\subsubsection*{B.3. Cas $x<0$}

Si $x<0$, alors :
\[
a_x=-\frac{k}{m}x>0.
\]
L'accélération est dirigée vers la droite, donc vers l'origine.

\begin{resultatbox}
\[
\boxed{x<0 \Longrightarrow a_x>0.}
\]
\end{resultatbox}

\subsubsection*{B.4. Accélération toujours dirigée vers l'origine}

Le signe de $a_x$ est toujours opposé au signe de $x$ :
\[
a_x=-\frac{k}{m}x.
\]
Ainsi, si le mobile est à droite de l'origine, l'accélération est vers la gauche.
S'il est à gauche de l'origine, l'accélération est vers la droite.

\begin{resultatbox}
\[
\boxed{\vec a \text{ est toujours dirigée vers } O.}
\]
\end{resultatbox}

\subsection{Partie C -- Étude qualitative du mouvement}

\subsubsection*{C.1. Après un lâcher vers la droite sans vitesse initiale}

On écarte le mobile vers la droite :
\[
x_0>0,
\]
puis on le lâche sans vitesse :
\[
v_0=0.
\]
Comme $x_0>0$, la force et l'accélération sont dirigées vers la gauche.
Le mobile commence donc à se mettre en mouvement vers l'origine.

\begin{resultatbox}
\[
\boxed{\text{Le mobile accélère vers l'origine dès le lâcher.}}
\]
\end{resultatbox}

\subsubsection*{C.2. Passage par l'origine}

À l'origine :
\[
x=0,
\]
donc :
\[
\vec F_r=\vec 0,
\qquad
\vec a=\vec 0.
\]
Mais la vitesse acquise n'est pas nulle : le mobile traverse l'origine par inertie.

\begin{attentionbox}
Force nulle ne signifie pas vitesse nulle. Cela signifie seulement accélération nulle à cet instant.
\end{attentionbox}

\subsubsection*{C.3. Pourquoi un mouvement d'aller-retour ?}

Une fois passé de l'autre côté, par exemple avec $x<0$, la force devient dirigée vers la droite.
Elle freine alors le mobile, puis le ramène vers l'origine.
Le même phénomène se répète de part et d'autre.

\begin{resultatbox}
\[
\boxed{\text{On s'attend à un mouvement oscillant autour de l'origine.}}
\]
\end{resultatbox}

\begin{remarquebox}
Ce modèle est celui d'un oscillateur harmonique sans frottements :
\[
a_x=-\frac{k}{m}x.
\]
C'est la même structure mathématique qu'un ressort idéal.
\end{remarquebox}

\subsection{Partie D -- Étude quantitative}

Données :
\[
m=\SI{0.400}{kg},\qquad k=\SI{5.0}{N.m^{-1}}.
\]

\subsubsection*{D.1. Force pour $x=\SI{0.10}{m}$}

\[
F_r=-kx=-5.0\times 0.10=-0.50\ \si{N}.
\]

\begin{resultatbox}
\[
\boxed{\vec F_r=-0.50\,\vec e_x\ \si{N}.}
\]
La valeur de la force est $\boxed{0.50\ \si{N}}$.
\end{resultatbox}

\subsubsection*{D.2. Accélération correspondante}

\[
a_x=\frac{F_x}{m}=\frac{-0.50}{0.400}=-1.25\ \si{m.s^{-2}}.
\]

\begin{resultatbox}
\[
\boxed{a_x=-1.25\ \si{m.s^{-2}}.}
\]
\end{resultatbox}

\subsubsection*{D.3. Cas $x=\SI{0.25}{m}$}

Force :
\[
F_r=-kx=-5.0\times 0.25=-1.25\ \si{N}.
\]
Accélération :
\[
a_x=\frac{-1.25}{0.400}=-3.125\ \si{m.s^{-2}}.
\]
Avec un nombre raisonnable de chiffres significatifs :
\[
a_x\simeq -3.13\ \si{m.s^{-2}}.
\]

\begin{resultatbox}
\[
\boxed{\vec F_r=-1.25\,\vec e_x\ \si{N}},
\qquad
\boxed{a_x\simeq -3.13\ \si{m.s^{-2}}.}
\]
\end{resultatbox}

\subsubsection*{D.4. Comparaison}

Lorsque $x$ passe de $\SI{0.10}{m}$ à $\SI{0.25}{m}$, il est multiplié par :
\[
\frac{0.25}{0.10}=2.5.
\]
Comme :
\[
a_x=-\frac{k}{m}x,
\]
l'accélération est proportionnelle à $x$. Sa valeur est donc aussi multipliée par $2.5$ :
\[
\frac{3.125}{1.25}=2.5.
\]

\begin{resultatbox}
\[
\boxed{\text{Plus le mobile est éloigné de l'origine, plus le rappel est intense.}}
\]
\end{resultatbox}

\subsection{Partie E -- Comparaison avec le cours}

\subsubsection*{E.1. Différence avec le poids}

Le poids vaut :
\[
\vec P=m\vec g.
\]
Près de la surface terrestre, il est vertical, dirigé vers le bas, et de valeur approximativement constante pour une masse donnée.

La force de recentrage vaut :
\[
\vec F_r=-kx\vec e_x.
\]
Elle est horizontale et dépend de la position $x$.

\begin{resultatbox}
\[
\boxed{\text{Le poids est vertical et presque constant ; la force de recentrage est horizontale et dépend de }x.}
\]
\end{resultatbox}

\subsubsection*{E.2. Différence avec une force de frottement}

Une force de frottement s'oppose généralement au mouvement, donc à la vitesse.
Par exemple :
\[
\vec f=-\lambda \vec v.
\]
La force de recentrage, elle, dépend de la position :
\[
\vec F_r=-kx\vec e_x.
\]
Elle ne s'oppose pas toujours à la vitesse. Lors du retour vers l'origine, elle peut être dans le même sens que la vitesse.

\begin{resultatbox}
\[
\boxed{\text{Le frottement dépend du mouvement ; le recentrage dépend de la position.}}
\]
\end{resultatbox}

\subsubsection*{E.3. Exemple réel}

Un exemple réel est la force de rappel exercée par un ressort idéal :
\[
\vec F=-k x\,\vec e_x,
\]
lorsque $x$ est l'allongement par rapport à la position d'équilibre.

Autres exemples acceptables :
\begin{itemize}
  \item pendule pour de petites oscillations ;
  \item membrane vibrante ;
  \item suspension de voiture modélisée par un ressort.
\end{itemize}

\subsection{Partie F -- Ajout de frottements}

On ajoute :
\[
\vec f=-\lambda \vec v,
\qquad \lambda>0.
\]

\subsubsection*{F.1. Effet qualitatif}

La force de frottement est opposée à la vitesse. Elle dissipe l'énergie mécanique.
Le mouvement d'aller-retour persiste au début, mais son amplitude diminue progressivement.

\begin{resultatbox}
\[
\boxed{\text{Les oscillations sont amorties.}}
\]
\end{resultatbox}

\subsubsection*{F.2. Le mobile finit-il par s'arrêter ? Où ?}

Oui, on peut s'attendre à ce que le mobile finisse par s'arrêter, car les frottements dissipent l'énergie.
La position finale stable est l'origine :
\[
x=0.
\]
En effet, à l'équilibre final :
\[
\vec v=\vec 0
\]
donc la force de frottement est nulle, et il faut aussi :
\[
\vec F_r=\vec 0.
\]
Or :
\[
\vec F_r=\vec 0 \Longleftrightarrow x=0.
\]

\begin{resultatbox}
\[
\boxed{\text{Le mobile finit au repos en }x=0.}
\]
\end{resultatbox}

\subsubsection*{F.3. Pourquoi le modèle est-il plus réaliste ?}

Dans la réalité, les frottements ne sont jamais exactement nuls : contacts, air, déformations, échauffement.
Un modèle sans frottements prévoit des oscillations éternelles, ce qui est idéal.
L'ajout d'un frottement permet de rendre compte de la perte d'énergie et de l'arrêt progressif du mobile.

\begin{resultatbox}
\[
\boxed{\text{Le modèle avec frottements est plus réaliste car il tient compte de la dissipation d'énergie.}}
\]
\end{resultatbox}

%========================================================
\newpage
\section{Bilan des méthodes à retenir}

\begin{methodebox}{Résoudre un problème de mécanique en Terminale, niveau approfondi}
\begin{enumerate}
  \item Définir clairement le système.
  \item Choisir un référentiel et préciser s'il est supposé galiléen.
  \item Faire un bilan complet des forces.
  \item Appliquer la deuxième loi de Newton :
  \[
  \sum \vec F_{\text{ext}}=m\vec a.
  \]
  \item Projeter sur des axes adaptés.
  \item Intégrer l'accélération pour obtenir la vitesse, puis la position.
  \item Utiliser les conditions initiales.
  \item Exploiter les équations obtenues.
  \item Vérifier les unités et interpréter physiquement.
\end{enumerate}
\end{methodebox}

\begin{methodebox}{Erreurs classiques}
\begin{itemize}
  \item Confondre vitesse et accélération.
  \item Dire qu'un mouvement uniforme n'est jamais accéléré.
  \item Oublier que l'accélération est vectorielle.
  \item Ajouter une force fictive de mouvement dans un référentiel galiléen.
  \item Oublier de convertir les km/h en m/s.
  \item Ne pas distinguer composante horizontale et composante verticale.
  \item Penser que force nulle implique vitesse nulle.
\end{itemize}
\end{methodebox}

\begin{center}
\Large\textbf{Fin du corrigé}
\end{center}

\end{document}
